
Automated code review pipelines increasingly rely on LLM-generated confidence signals to prioritize or filter code changes. If the probability a model assigns to its own output being correct (P(True)) degrades systematically on problems the model cannot solve, practitioners can compensate with difficulty-aware thresholds. If the degradation pattern is absent or inverted, a single threshold calibrated on average performance will misfire on subsets of the problem distribution.

This study was designed to test whether harder code problems cause greater calibration error than easy problems, measured as DELTA\_ECE = ECE(hard) - ECE(easy) > 0. The pre-registered prediction was that this positive DELTA\_ECE would hold in at least 2 of 3 tested model families. This prediction was refuted. The actual finding is that the direction of DELTA\_ECE depends on model architecture and training regime.

The surface problem is well-established: modern neural networks, including LLMs, are poorly calibrated overall \citep{Guo2017Calibration}. LLMs exhibit P(True) confidence signals that do not reliably correspond to actual correctness rates \citep{Kadavath2022Language}. In code generation evaluation, pass@k metrics measure output quality \citep{Chen2021HumanEval} and EvalPlus augmented tests provide reliable correctness oracles \citep{Liu2023EvalPlus}, but neither line of work examines whether calibration quality varies with problem difficulty.

The deeper problem is that global calibration measures obscure difficulty-conditional structure. If ECE is systematically higher on hard problems than easy ones, a model's confidence is less reliable exactly when it matters most. Conversely, if ECE is higher on easy problems (as observed for CodeLlama), the model is overconfident on problems it should handle well.

No existing work applies difficulty-stratified calibration analysis to LLM code verification. P(True) calibration studies examine capability scaling across model sizes, not task-difficulty conditioning. Code evaluation benchmarks measure pass@k but not calibration quality. The combination of P(True), ECE, and self-contained difficulty bootstrapped from the experiment's own pass@1 distribution has not been studied.

This work makes the following contributions:

\begin{enumerate}
\item A validated self-contained calibration methodology: a four-step pipeline (k=5 solution generation, tier stratification, P(True) logprob extraction, ECE computation) achieving perfect coverage (1.0000) for all model-benchmark combinations with non-degenerate confidence distributions (std 0.062-0.078).
\end{enumerate}

\begin{enumerate}
\item An architecture-dependent calibration finding: three architectures representing distinct training regimes produce qualitatively different DELTA\_ECE patterns. The pre-registered universal prediction (positive DELTA\_ECE in at least 2/3 architectures) is refuted, but the architecture-stratified result is more informative.
\end{enumerate}

\begin{enumerate}
\item Evidence consistent with a training-data composition mechanism: temperature scaling fails to reverse DELTA\_ECE direction for CodeLlama (post-scaling DELTA\_ECE worsens to -0.810) while DeepSeek's positive DELTA\_ECE persists (+0.073). These observations are exploratory (N=1 per architecture category).
\end{enumerate}

\begin{enumerate}
\item A reusable difficulty-calibration infrastructure validated on EvalPlus with k=5 self-contained difficulty bootstrap, P(True) extraction, and M=15-bin ECE with bootstrap confidence intervals.
\end{enumerate}
